\subsection{DIMENSIONS OF SIMPLE g-MODULES}

If \(\mu \in \mathfrak{h}^*\) , put \(\mathrm{J}(e^{\mu}) = \sum_{w\in \mathrm{W}}\varepsilon (w)e^{w\mu}\) , cf.~Chap. VI, §3, no. 3.

THEOREM 2. Let E be a finite dimensional simple g-module, \(\lambda\) its highest weight and \((\cdot|\cdot)\) a W-invariant non-degenerate positive symmetric bilinear form on \(h^{*}\) . Then:

\[
\dim \mathrm{E} = \prod_ {\alpha \in \mathrm{R} _ {+}} \frac {\langle \lambda + \rho , H _ {\alpha} \rangle}{\langle \rho , H _ {\alpha} \rangle} = \prod_ {\alpha \in \mathrm{R} _ {+}} \left(1 + \frac {(\lambda | \alpha)}{(\rho | \alpha)}\right).
\]

Let \(\mathrm{T}\) be an indeterminate. For all \(\nu \in \mathrm{P}\) , denote by \(f_{\nu}\) the homomorphism from \(\mathbf{Z}[\mathrm{P}]\) to \(\mathbf{R}[[\mathrm{T}]]\) that takes \(e^{\mu}\) to \(e^{(\nu |\mu)\mathrm{T}}\) for all \(\mu \in \mathrm{P}\) . Then \(\dim \mathrm{E}\) is the constant term of the series \(f_{\nu}(\operatorname{ch}\mathrm{E})\) .

For all \(\mu, \nu \in \mathrm{P}\) , we have

\[
\begin{array}{c} f _ {\nu} (\mathrm{J} (e ^ {\mu})) = \sum_ {w \in \mathrm{W}} \varepsilon (w) e ^ {(\nu | w \mu) \mathrm{T}} \\ = \sum_ {w \in \mathrm{W}} \varepsilon (w) e ^ {(w ^ {- 1} \nu | \mu) \mathrm{T}} = f _ {\mu} (\mathrm{J} (e ^ {\nu})). \end{array}
\]

In particular, in view of Chap. VI, §3, no. 3, formula (3),

\[
f _ {\rho} (\mathrm{J} (e ^ {\mu})) = f _ {\mu} (\mathrm{J} (e ^ {\rho})) = e ^ {(\mu | \rho) \mathrm{T}} \prod_ {\alpha \in \mathrm{R} _ {+}} (1 - e ^ {- (\mu | \alpha) \mathrm{T}}).
\]

Hence, setting \(\mathrm{Card}(\mathrm{R}_{+})=\mathrm{N}\) ,

\[
f _ {\rho} (\mathrm{J} (e ^ {\mu})) \equiv \mathrm{T} ^ {\mathrm{N}} \prod_ {\alpha \in \mathrm{R}} (\mu | \alpha) \pmod {\mathrm{T} ^ {\mathrm{N} + 1} \mathbf {R} [ [ \mathrm{T} ] ]}.
\]

The equality \(\mathrm{J}(e^{\lambda +\rho}) = \mathrm{ch}(\mathrm{E}).\mathrm{J}(e^{\rho})\) (Th. 1) thus implies that

\[
\mathrm{T} ^ {\mathrm{N}} \prod_ {\alpha \in \mathrm{R} _ {+}} (\lambda + \rho | \alpha) \equiv f _ {\rho} (\operatorname{ch} \mathrm{E}). \mathrm{T} ^ {\mathrm{N}} \prod_ {\alpha \in \mathrm{R} _ {+}} (\rho | \alpha) \pmod {\mathrm{T} ^ {\mathrm{N} + 1} \mathbf {R} [ [ \mathrm{T} ] ]}
\]

SO

\[
\dim \mathrm{E} = \left(\prod_ {\alpha \in \mathrm{R} _ {+}} (\lambda + \rho | \alpha)\right) / \left(\prod_ {\alpha \in \mathrm{R} _ {+}} (\rho | \alpha)\right) = \prod_ {\alpha \in \mathrm{R} _ {+}} \left(1 + \frac {(\lambda | \alpha)}{(\rho | \alpha)}\right).
\]

Now, if \(\alpha \in \mathrm{R}_{+}\) , \(\alpha\) can be identified with an element of \(\mathfrak{h}_{\mathbf{R}}\) proportional to \(H_{\alpha}\) , so

\[
(\lambda + \rho | \alpha) / (\rho | \alpha) = \langle \lambda + \rho , H _ {\alpha} \rangle / \langle \rho , H _ {\alpha} \rangle .
\]

Examples. 1) In the Example of no. 1, we find that

\[
\dim \mathrm{V} (m) = \left(\frac {m}{2} \alpha + \frac {\alpha}{2}\right) (H _ {\alpha}) / \frac {\alpha}{2} (H _ {\alpha}) = m + 1,
\]

which we knew in §1.

\begin{enumerate}
\def\labelenumi{\arabic{enumi})}
\setcounter{enumi}{1}
\tightlist
\item
  Take g to be the splittable simple Lie algebra of type \(G_{2}\) and adopt the notations of Chap. VI, Plate IX. Give \(h_{R}^{*}\) the W-invariant positive symmetric form \((\cdot|\cdot)\) such that \((\alpha_{1}|\alpha_{1})=1\) . Then \(\rho=\varpi_{1}+\varpi_{2}\) and
\end{enumerate}

\[
\begin{array}{c} (\varpi_ {1} | \alpha_ {1}) = \frac {1}{2}, (\varpi_ {1} | \alpha_ {2}) = 0, (\varpi_ {1} | \alpha_ {2} + \alpha_ {1}) = \frac {1}{2}, \\ (\varpi_ {1} | \alpha_ {2} + 2 \alpha_ {1}) = 1, (\varpi_ {1} | \alpha_ {2} + 3 \alpha_ {1}) = \frac {3}{2}, (\varpi_ {1} | 2 \alpha_ {2} + 3 \alpha_ {1}) = \frac {3}{2}, \\ (\varpi_ {2} | \alpha_ {1}) = 0, (\varpi_ {2} | \alpha_ {2}) = \frac {3}{2}, (\varpi_ {2} | \alpha_ {2} + \alpha_ {1}) = \frac {3}{2}, \\ (\varpi_ {2} | \alpha_ {2} + 2 \alpha_ {1}) = \frac {3}{2}, (\varpi_ {2} | \alpha_ {2} + 3 \alpha_ {1}) = \frac {3}{2}, (\varpi_ {2} | 2 \alpha_ {2} + 3 \alpha_ {1}) = 3. \end{array}
\]

Thus, if \(n_1, n_2\) are integers \(\geq 0\) , the dimension of the simple representation of highest weight \(n_1\varpi_1 + n_2\varpi_2\) is

\[
\begin{array}{c} \left(1 + \frac {n _ {1} / 2}{\frac {1}{2}}\right) \left(1 + \frac {3 n _ {2} / 2}{\frac {3}{2}}\right) \left(1 + \frac {n _ {1} / 2 + 3 n _ {2} / 2}{\frac {1}{2} + \frac {3}{2}}\right) \left(1 + \frac {n _ {1} + 3 n _ {2} / 2}{1 + \frac {3}{2}}\right) \\ \times \left(1 + \frac {3 n _ {1} / 2 + 3 n _ {2} / 2}{\frac {3}{2} + \frac {3}{2}}\right) \left(1 + \frac {3 n _ {1} / 2 + 3 n _ {2}}{\frac {3}{2} + 3}\right) \\ = (1 + n _ {1}) (1 + n _ {2}) \left(1 + \frac {n _ {1} + 3 n _ {2}}{4}\right) \left(1 + \frac {2 n _ {1} + 3 n _ {2}}{5}\right) \left(1 + \frac {n _ {1} + n _ {2}}{2}\right) \\ \times \left(1 + \frac {n _ {1} + 2 n _ {2}}{3}\right) \\ = \frac {(1 + n _ {1}) (1 + n _ {2}) (2 + n _ {1} + n _ {2}) (3 + n _ {1} + 2 n _ {2}) (4 + n _ {1} + 3 n _ {2}) (5 + 2 n _ {1} + 3 n _ {2})}{5 !}. \end{array}
\]

In particular, the fundamental representation of highest weight \(\varpi_{1}\) (resp. \(\varpi_{2}\) ) is of dimension 7 (resp. 14).

